%% =====================================================================
%%  core/config/margenes.tex  --  Geometria de pagina (A4 oficial)
%% ---------------------------------------------------------------------
%%  Replica el machote del despacho: A4 con margenes de 1 pulgada
%%  (2.54 cm) a los cuatro lados. El membrete vive en la zona de
%%  encabezado (fancyhdr), por eso headheight es generoso: la caja del
%%  nombre + logos + lema ocupan ~3 cm.
%% =====================================================================

\RequirePackage[
  a4paper,
  includehead,         % el membrete cuenta dentro del margen superior
  % Zona superior ajustada al ALTO REAL del membrete (Edison
  % 2026-09-07: «la parte superior ocupa demasiado espacio»). El
  % membrete mide ~1.98 cm (logos + caja del nombre + lema); con
  % headheight=2.60 sobraban ~0.5 cm de blanco ARRIBA -- la minipage
  % va con [b], así que la holgura de la caja cae encima del arte.
  % Del borde al primer texto: 44.2 mm -> ~33 mm. Si se cambia el
  % membrete por uno más alto, subir headheight (fancyhdr avisa en el
  % log: «\headheight is too small»).
  top=1.00cm,          % borde del papel -> inicio del membrete
  % 2.15 cm = 61.2pt. fancyhdr exige >= 60.59pt para el membrete
  % estandar (oficio/memorando/informe); con 2.10 avisaba en el log.
  headheight=2.15cm,   % alto reservado al membrete (minimo real: 60.59pt)
  headsep=0.35cm,      % membrete <-> cuerpo
  left=2.54cm,
  right=2.54cm,
  % Banda del pie de 1.50 cm (antes 1.10) para que el SELLO de consigna
  % quepa DENTRO de ella, bajo las notas al pie y pegado al filete, sin
  % pisar el texto (Edison 2026-09-07). El area de texto solo cede 0.16 cm.
  bottom=2.70cm,
  footskip=1.50cm,
]{geometry}
